\documentclass[preprint, 3p,
authoryear]{elsarticle} %review=doublespace preprint=single 5p=2 column
%%% Begin My package additions %%%%%%%%%%%%%%%%%%%

\usepackage[hyphens]{url}

  \journal{CERC 2026?} % Sets Journal name

\usepackage{graphicx}
%%%%%%%%%%%%%%%% end my additions to header

\usepackage[T1]{fontenc}
\usepackage{lmodern}
\usepackage{amssymb,amsmath}
% TODO: Currently lineno needs to be loaded after amsmath because of conflict
% https://github.com/latex-lineno/lineno/issues/5
\usepackage{lineno} % add
\usepackage{ifxetex,ifluatex}
\usepackage{fixltx2e} % provides \textsubscript
% use upquote if available, for straight quotes in verbatim environments
\IfFileExists{upquote.sty}{\usepackage{upquote}}{}
\ifnum 0\ifxetex 1\fi\ifluatex 1\fi=0 % if pdftex
  \usepackage[utf8]{inputenc}
\else % if luatex or xelatex
  \usepackage{fontspec}
  \ifxetex
    \usepackage{xltxtra,xunicode}
  \fi
  \defaultfontfeatures{Mapping=tex-text,Scale=MatchLowercase}
  \newcommand{\euro}{€}
\fi
% use microtype if available
\IfFileExists{microtype.sty}{\usepackage{microtype}}{}
\usepackage[]{natbib}
\bibliographystyle{plainnat}

\ifxetex
  \usepackage[setpagesize=false, % page size defined by xetex
              unicode=false, % unicode breaks when used with xetex
              xetex]{hyperref}
\else
  \usepackage[unicode=true]{hyperref}
\fi
\hypersetup{breaklinks=true,
            bookmarks=true,
            pdfauthor={},
            pdftitle={Tactical urbanism for mobility hubs? DIY mapping with mobilityHubTools},
            colorlinks=false,
            urlcolor=blue,
            linkcolor=magenta,
            pdfborder={0 0 0}}

\setcounter{secnumdepth}{5}
% Pandoc toggle for numbering sections (defaults to be off)


% tightlist command for lists without linebreak
\providecommand{\tightlist}{%
  \setlength{\itemsep}{0pt}\setlength{\parskip}{0pt}}

% From pandoc table feature
\usepackage{longtable,booktabs,array}
\usepackage{calc} % for calculating minipage widths
% Correct order of tables after \paragraph or \subparagraph
\usepackage{etoolbox}
\makeatletter
\patchcmd\longtable{\par}{\if@noskipsec\mbox{}\fi\par}{}{}
\makeatother
% Allow footnotes in longtable head/foot
\IfFileExists{footnotehyper.sty}{\usepackage{footnotehyper}}{\usepackage{footnote}}
\makesavenoteenv{longtable}



\usepackage{subfig}
\usepackage{booktabs}
\usepackage{longtable}
\usepackage{array}
\usepackage{multirow}
\usepackage{wrapfig}
\usepackage{float}
\usepackage{colortbl}
\usepackage{pdflscape}
\usepackage{tabu}
\usepackage{threeparttable}
\usepackage{threeparttablex}
\usepackage[normalem]{ulem}
\usepackage{makecell}
\usepackage{xcolor}



\begin{document}


\begin{frontmatter}

  \title{Tactical urbanism for mobility hubs? DIY mapping with
mobilityHubTools}
    \author[]{James Reynolds%
  \corref{cor1}%
  \fnref{1}}
   \ead{julianj.reynolds@atu.ie} 
    \author[]{Brian McCann%
  %
  \fnref{1}}
   \ead{brian.mccann@atu.ie} 
    \author[]{Noreen Armstrong%
  %
  \fnref{2}}
   \ead{noreen.armstrong@atu.ie} 
    \author[]{Sarbast Molsem%
  %
  \fnref{3}}
   \ead{molsem@tcd.ie} 
    \author[]{Agnieszka Stefaniec%
  %
  \fnref{4}}
   \ead{a.stefaniec@soton.ac.uk} 
    \author[]{Brian Caulfield%
  %
  \fnref{3}}
   \ead{brian.caulfield@tcd.ie} 
    \author[]{Amaya Vega%
  %
  \fnref{2}}
   \ead{amaya.Vega@atu.ie} 
      \cortext[cor1]{Corresponding author}
    \fntext[1]{Atlantic Technological University, Sligo campus, Ireland}
    \fntext[2]{Atlantic Technological University, Galway City campus,
Ireland}
    \fntext[3]{Centre for Transport Research, Dept of Civil, Structural
and Environmental Engineering, Trinity College Dublin, Ireland}
    \fntext[4]{Department of Decision Analytics and Risk at Southampton
Business School, University of Southampton, UK}
  
  \begin{abstract}
  
  \end{abstract}
  
 \end{frontmatter}

\section{Introduction}\label{introduction}

`Mobility hub' is a relatively new term for locations where a range of
transport modes, especially transit and shared micro-mobility, and
sometimes other services and facilities are offered. These might provide
a range of non-private-or-petrol-driven-car alternatives that might help
towards more sustainable travel behaviours and transport system
outcomes. Nearby shared electric cars, bikes and cargo-bikes, e-bikes
and e-cargo bikes, or other vehicles might help in shifting towards more
walking, cycling and transit-based travel through providing options for
those, more occasional, trips that are otherwise difficult without a
car, (e.g.~travelling with luggage, longer distances or outside typical
service times). A nearby mobility hub might allow households to forgo
purchase of a(n additional) car, or transition to electrically-powered
modes without the upfront expense. Some mobility hubs, such as those in
the Greater Toronto and Hamilton Area (GTHA) in Canada, are defined as
such by a regional transport authority, with implications for landuse
planning and other institutional decision-making
\citep{Salsberg:2010aa, Engel-Yan:2012aa}. Others are created afresh,
through some agency implementing facilities for a range of modes and/or
other offerings in close proximity to each other and with branding as a
`mobility hub' or similar.

In some cases it may be that a range of offerings are already in place,
especially where the private sector has already implemented car-,
scooter-, bike- or other -share services. Converting such places into a
`mobility hub' may be within reach with just a mobility-hub-related
name, a press-release or website, and perhaps some minimal physical
interventions such as signage and advertising. However, tools for
transforming such locations into `mobility hubs' do not yet appear to be
readily available. Instead, individual agencies implementing mobility
hubs appear to have been largely created signage, advertising and
similar elements afresh and bespoke by on a site-by-site basis, albeit
that some hubs created in groups appear to have similar branding and
designs. Such creation might involve graphic design, mapping or similar
and so limit the opportunities for trials, bottom-up implementation or
other low-cost interventions by practitioners or (potentially through
tactical urbanism and/or grassroots activism) sustainability-focused
advocates or other actors.

This paper reports findings from case and design-based research seeking
to develop new tools to support low-cost mobility hub implementation.
The research adopts the Double Diamond Design process and iterates
through alternatively diverging and coverging phases of Discovery,
Definition, Development and Delivery (c.f.
\cite{UK-Design-Council:0aa}). Exploratory field visits to a range of
case study mobility hubs help inform the Discovery and Development
phases, while a new package of software tools to support mapping and the
generation of informational signage and other physical interventions
forms the key research outcome from the Delivery phase. This software
might allow practitioners or others to create a temporary mobility hub
on a trial basis, or as the first step towards designing and
implementing more permanent information and mapping. The rest of this
paper is structured as follows: the next section outlines the context
for this research, including brief discussion of mobility hubs and
tactical urbanism. The research methodology is then described, followed
by presentation of the results, which are then briefly discussed in
conjunction with a conclusion identifying directions for future
research.

\section{Research context}\label{research-context}

Mobility hubs are an emerging approach where shared and micro-mobility
offerings, transit or other offerings are linked to provide a range of
options that support more sustainable travel. Early use of the term
occured in Toronto, Canada, as part of landuse and transport planning
undertaken by Metrolinx and other governmental institutions across the
surrounding region \citep{Salsberg:2010aa, Engel-Yan:2012aa}. In ``The
Big Move'' \citep{Metrolinx:2008aa} various transit stations and major
interchanges were defined as ``mobility hubs'' for the purposes of
supporting regional travel, modal interchange, and transitions to more
transit supportive landuse patterns.

The plethora of independent private-sector transport offerings now
available, including car-share and dockless bike/scooter systems, may
also have helped towards development of the mobility hub concept.
Although driven largely by start-up, technology and communications, and
related sectors, various governmental agencies have sought to both:
regulate such offerings; and make them more than the sum of their parts.
Defining dockless parking locations, allocating dedicated carshare
parking spaces, installation of branding and informational signage and
similar efforts may help to provide geo-spatial context and
collectivisation. As well, there have been some emerging problems,
including with respect to inappropriate parking of dockless
bikes/scooters (i.e.~blocking footpaths), and risks and environmental
impacts of recharging models that rely on batteries being collected and
then charged elsewhere, often without the appropriate facilities to deal
with electrical loads or battery fires
\citep{Meigs:2023aa, Zhang:2022aa}. Shared electric mobility hubs, such
as those currently being trialled in Ireland as part of the ROBUST
project \citep{ROBUST:2025aa}, may provide solutions to such problems by
shifting the recharging of vehicles to be at a hub itself where
appropriate parking, electrical connections and other equipment
(importantly for fires involving lithium batteries) are outside and
typically unenclosed.

Other challenges for implementers and communities, however, may relate
to the difficulty of developing mobility hubs, including the planning
processes and funding acquisition necessary to build new facilities. Of
course, other than the name there may be nothing new about mobility hubs
- railway stations have existing for over a century. Often containing
waiting areas and information, cafes, parking lots, interchange with
other modes such as buses, and many other facilities the 19th, 20th and
current iterations of railway stations would appear to have largely
already `solved', many of the problems related to how to support
mobility at a network hub (albeit that these solutions might not always
be fully implemented). However, stations tend to be located on land
owned or otherwise controlled by railway companies or agencies. As such,
facilities, signage, branding, information displays and most other
aspects of the station environment might be largely under the control of
a single institution, or only a small number of closely-related
agencies. Mobility hubs, in contrast, might tend to be located in public
road reserves, within parking facilities or otherwise where there may be
less clear authority and power to make changes without widespread
public, institutional and government approvals or buy-in. These may also
often require outside expertise and/or involve top-down policy-making in
which those making decisions might be somewhat removed from the real
mobility needs of individual communities, or who may adopt `one size
fits all' approaches lacking responsiveness to local context.

Research in public policy analysis has long been examining on such
non-rational factors and how they might influence decision-making, but
these tend to be overlooked in transport policy research
\citep{Marsden:2017aa}. Instutionalism, political approaches,
incrementalism and other public policy models, however, are clearly
applicable to transport systems (see for example
\citet{Reynolds:2017ah}). Important within these are the top-down,
bottom-up and hybrid theories of implementation, which emphasise the
influence that street-level bureaucrats can wield to obtain outcomes
that respond to their local needs and context, including sometimes
outcomes that might be in opposition to central policy directions.
Further exploration of this topic may be beyond the scope of this paper,
but the extensive body of research related to the Advocacy Coalition
Framework (including
\citet{Sabatier:1986aa,Sabatier:1987aa,Sabatier:1988aa,Sabatier:1993aa,Weible:2009aa,Weible:2007aa})
provides insights and theory.

Perhaps of more relevance, however, is the field of tactical urbanism,
which adopts many such bottom-up and grassroots approaches to defining
and re-imaging public spaces, especially those spaces within road
reserves and related to transport systems. Tactical urbanism can involve
making changes to public environments, often unsanctioned, that
sometimes question what urban space is for or the legitimacy of how it
is allocated. ``PARKing Day'' provides an example, having started as
lunchtime take over of an on-street car parking space by an architecture
firm, who rolled out some green grass and a park bench, paid the parking
meter, and in doing so questioned the validity of using such spaces for
storing vehicles instead of as a parklet for people. This has since
grown into an annual event seeking, like the World Naked Bike Ride and
Critical Mass, to demonstrate the possibility of change and question the
legitimacy of the (car-dominated) status quo
\citep{Lydon:2012aa,Davidson:2013aa,Lydon:2015aa,Blumgart:2016ab, Ride:2009aa, Morhayim:2012aa, Longhurst:2015aa}.
Similar `pop-up' or protest-based approaches have helped legitimise
bicycle lanes, delegitimise car-use of bus lanes, and are sometimes used
as more mainstream ways (amongst engineers and planners) to help build
public and institutional support for implementation, including for the
installation of transit lanes
\citep{Fucoloro:2013aa,News:2019ab, TransitCenter:2018ab, Transportation-Studies:2019aa}.
These build on or have similarities to various earlier approaches
including street-reclaiming
\citep{Dobson:2023ae,Engwicht:1999aa,Fotel:2009aa,Institute:2015aa}, the
invention of the first `woornef' \citep{Guttenberg:1982aa, Lydon:2015aa}
and Local Area Traffic Management \citep{Damen:2017aa}, as well as being
evident in the make up of current directions towards `Movement and
Place', `Complete Streets' and Network Operation Planning approaches
\citep{Delbosc:2018aa}.

Despite the many names, variations and focii of these and other
interventional and decision-making approaches, the call from
\citet{Marsden:2017aa} to ``pay greater attention to context, politics,
power, resources and legitimacy'' retains relevance. Participating in
unsanctioned activities, especially for minorities or others who might
attract higher levels of attention from authorities, can be a risk that
might limit engagement in tactical urbanism or similar activities
\citep{Douglas:2018aa}. Intersections between illegality, illegitimacy
and the ill-enforced similarly exist in the context of street art and
its place on a continuum that includes commissioned murals,
counter-culturism and the \emph{avant-garde}, graffiti and vandalism.
\citet{Killen:2017aa,Killen:2021aa} explore these topics in the context
of trains, but of relevance here, however, are questions relating to who
is allowed to post a public notice, define a `mobility hub' or otherwise
share knowledge about transport options, local facilities and services,
and connecting modes.

Engineers, planners, operational managers, communications teams and
others may have less need to engage in tactical-urbanist-type
implementation given that they might already have legitimate authority
to make changes to the built environment, at least within the scope of
their employment and remit. However, pop-ups, trials, bottom-up and
incremental approaches to implementation may still be of use, or even
necessary, for such actors when seeking legitimacy for new, unusual or
just controversial interventions.
\citet{Reynolds:2020aa,Reynolds:2021ab} define these as `pragmatic
strategies for legitimation \emph{through} implementation' based on case
research of such approaches helping in the implementation of Bus Rapid
Transit in Curitiba and other forms of transit priority in Zurich,
Melbourne and Toronto. Other strategies include building legitimacy
\emph{before} implementation using technical evaluation, vision-based
planning and/or public decision-making processes; or avoiding impacts on
those who might seek to delegitimise change.

Basic signage, linemarking or the public display of travel information
would appear likely to be relatively easy for many practitioners to
implement within their own institutional setting. What may be lacking,
however, are tools and standarised approaches to mobility hub
information provision for quick and relatively easy adaptation to local
contexts. Likewise, for activists, practitioners wishing to affect
change outside the remit of their specific employment suituation, or the
general public there do not appear to be tools that might support
tactical urbanism efforts related to mobility hubs. This might be
particularly relevant for the `creation' of a mobility hub out of
already existing and geographically-proximate facilitates, which might
not be immediately obvious to travellers but something informally
well-known and readily apparent to those with local knowledge of the
geography and community. Leaflets, posters for shop windows and
community noticeboards, or even resources for a tactical urbanist to
launch a `guerilla marketing' campaign of posters, Do-It-Yourself (DIY)
signage and websites, might be all that is needed to facilitiate
bottom-up action and empower communities to define their own hubs and
advertise their pre-existing mobility services and related facilities.

What appears to be lacking in the research literature, or in practice,
are the tools to quickly and easily produce the sorts of maps, analysis
and interpretations that might support the definition of a place as a
hub and the communciation of its existing and potential mobility and
other offerings. Mapping, signage and similar interventions in both the
physical and virtual might usually tend to be produced on a bespoke (ie.
location-by-location) basis by specialist geographers and graphic
designers, and involve commercially-controlled datasets or platforms.
With the advent of crowd-sourcing, open data and the widespread
availability of open source software, however, there is the possibilty
of developing location-bespoke, but generalised, mapping and similar
tools useful for practitioners or others wishing to `mobility hub-erise'
some part of the urban environment. The methods used in this research to
firstly define what such tools might do, and then to deliver a software
package that enables others to more easily develop their own mobility
hub mapping and informational materials, which might provide an initial
version for future refinement, are therefore descrbed in the following.

\section{Methodology}\label{methodology}

\citet{Coxon:2018aa,Coxon:2021ab} outline various industrial design
approaches and tools, as well as discussing their relevance and
potential applicability to transport systems. Notable amongst these is
the Double Diamond, which involves iteration across four phases
alternating between: diverging (1) Discovery and then converging (2)
Definition phases; followed by diverging (3) Development and convergence
to (4) Delivery of a design output \citep{UK-Design-Council:0aa}. This
study adopts this approach together with case studies of existing
mobility hubs.

Case studies can be especially effective when theory development is at
an early stage and when there is a need to explore phenomena within
real-world contexts, as is the situation here. A challenge, however, is
the need to be seeking generalisable findings despite researching at
depth within the case contexts. Such depth is possible through including
only a small number of cases, typically up to five to ten, yet leads to
the suggestion that researchers should explictly respond to this
``duality criterion''
\citep{Benbasat:1987aa,Eisenhardt:1989aa,Dyer:1991aa,Cavaye:1996aa,Darke:1998aa,Meredith:1998aa,Stuart:2002aa,Voss:2002aa,Denscombe:2007aa,Eisenhardt:2007aa,Yin:2009ab,Ketokivi:2014aa,Yin:2014ab,Yin:2018ab}.
Here the need to be seeking generalisable findings is addressed directly
by the form taken by outputs, namely: the research produces a software
of package tools for mobility hubs in any location, not just those
selected as cases.

Hubs were visited in London\footnote{South Woodford and High Street
  Wanstead; and in Camden: Charrington Street, Goldington Crescent,
  Werrington Street, Cranlegith Street, Brill Place, Doric Way and
  Chalton Street.}, Bristol\footnote{Abbey Wood, UWE Frenchay, Horfield
  Library and Gainsborough Square} and Brimingham\footnote{Huntingtree
  Park and Andrew Road}. Case selection for this study was somewhat
constrained by logistics of travel, meaning that visited sites were in
the south of England. However, and in line with the recommendations of
the methodological literature, the cases selected included leading and
particularly novel examples, those that are more representative, and
sufficient variation to provide for cross-case comparison. Cases were
selected from a larger sample available at
https://www.como.org.uk/mobility-hubs/built-and-planned-hubs.

The findings from site visits where synthesised through cross-case
comparison to develop the design problem Definition. The Development was
guided by standardised approaches to package development for the R
language, using \citet{wickham2023r} as a key reference. R is chosen
here because: it is open source and freely available, therefore
supporting goals of developing an accessible suite of tools; it has an
extensive set of packages for geospatial analysis and mapping, including
\emph{sf} \citep{R-sf}, and the \emph{osmdata} \citep{R-osmdata}
packages; and there is a large community of users who might help support
those with less experience in troubleshooting, but who are also
responsible for the development and maintenance of a wide range of
packages, including: the \emph{tidyverse} \citep{tidyverse2019} and many
others used in code development for this research.

\section{Results}\label{results}

\subsection{Discovery}\label{discovery}

\begin{figure}

{\centering \includegraphics[width=1\linewidth]{graphics/andrew_road} 

}

\caption{Andrew Road hub, selected images. Source: authors.}\label{fig:andrew_road_images}
\end{figure}

\begin{figure}

{\centering \includegraphics[width=1\linewidth]{graphics/bristol_03} 

}

\caption{Bristol hubs, selected images. Source: authors.}\label{fig:bristol_images_panels}
\end{figure}

Information panels, maps and other guides to facitilies and services
were present at some hubs, but notably not at those in Camden, South
Woodford or High Street Wanstead. These locations were all in London,
however, and generally close to already existing maps and other
information at tube stations, bus stops or otherwise provided, which may
have lessened the need for mobility-hub-specific signage and mapping.
Figures \ref{fig:andrew_road_images} and \ref{fig:bristol_images_panels}
show selected images of informational panels and signage found at the
hubs in Birmingham and Bristol, which shared some common design elements
such as including both local and area-wide maps, icongraphic
representation of facilities, and directional information. Notably, the
Bristol hubs included green cubes mounted overhead hub facilities, which
appear likely to assist in finding hub elements in busy urban
environments. Mapping and informational material, however, often
appeared to include some that are not specifically `part' of a hub
(e.g.~toliets within anouther facilities, shops, cafes). This highlights
a potential question for implementers: what services and facilities to
include as `part of the hub', versus those that are just `nearby'? In
the case of the Transport for West Midlands hub at Andrew Road (Figure
\ref{fig:andrew_road_images}) directional signage to, and information
about, bus services at the nearby Halesowen Bus Station (approx 50m
walking distance) is provided, but the hub has a different
location-based name and branding identity.

Some of the hubs included linemarked areas for shared e-bike, e-scooter
or e-cargobike parking, bicycle parking and secure parking lockers,
although the offered facilities and services varied by location. A lack
of shade, weather protection and drinking water was also evident at most
of the visited hubs, even though many of the parklet-type facilities
included steel framing that might have accommodated open-sided roofing
or similar. This might indicate an emphasis on form over function in
some design elements, whereas one of the hubs in Bristol perhaps
provides an example of additional functionality - a containerised
``Library of Things'' - that appears to have been added later into the
hub area, onto top of areas linemarked and sign-posted as for electric
scooter hire. This may provide an indication of the possibilities of
using modular components to expand, re-purpose or adjust mobility hubs
as situations change (consistent with \citet{Stuart:2002aa}), but a need
for signs and maps to be easily adjustable to match, as updating such
elements might relatively challenging for practitioners, location
managers or others once the original graphic designers or implementers
have moved on\footnote{See also \citet{Reynolds:2017ah} discussion of
  `The Garbage Can Model'.}.

Overall, the site visits undertaken during the `Discovery' phase
highlight the following key issues: design elements; included facilities
and hub identity; and the ability for supporting material to be adjusted
to changing circumstances. Design elements relate to the desirability of
functional and clearly-related hub components, and appears interlinked
clarifying what types of nearby facilities should be `part' of the hub
in terms of branding and information\footnote{i.e.~Which elements
  recieve `green cube' markers in the real world,\\
  versus which are just an icon on a map?}. Identity, likewise, appears
to interlinked with hub naming - for example, the Andrew Road hub is
named after a relatively less important connector road rather than
consolidated to be part of the Halesowen bus station.\\
As well the site visits perhaps also highlighted the importance of
maintenance, activity\footnote{c.f. \citet{Jacobs:1961aa} ``eye on the
  street''} and updates to match changing circumstances that might give
a hub a sense of `place'. Hubs in Bristol perhaps provide an example
with litter, non-mobility-related signage, space requisitioned for other
purposes, and some sites advertising dockless shared vehicles but
without these present in sufficient numbers (or at all at the time of
visiting Gainsborough Square) all providing indications of potential
post-implementation challenges related to mobility hubs.

\subsection{Definition}\label{definition}

Findings from the sites visits suggest a wide variety of modes, services
and facilities might be provided as part of the hub itself, as well as
further services and facilities identified as being `nearby' as an
extension of the hub offer. This might suggest adopting an open-ended
and expansionist approach to mobility hub definition.

The consistency of branding at the Bristol hubs, especially with respect
to the iconographic green cubes drawing attention to hub elements, might
highlight the importance of proximity and line-of-sight in defining a
mobility hub. Although mapping at most hubs provided information about
services and facilities within up to around 15 minutes walking distance,
hub elements appeared to generally be located within relatively confined
areas, generally spanning less than 50 metres. The Andrew Road hub
provides an example towards the upper end, with the entirety of the hub
and the nearby Halesown Bus Station encompassed within a radius of
approximately 70 metres, and spanning slightly less than 200 metres if
traversed by foot. In contrast, the South Woodford hub takes up
approximately 15 metres of kerb space and is less than 100 metres from
the nearest tube station, while some in Camden use space equivalent to
only a few car parking spots.

Definitions of the facilities, services and modes necessary to become a
mobility hub appear likewise to be open-ended and inclusive, rather than
some sort of minimum checklist that all sites must have\footnote{CoMoUK's
  pop-up mobility hub offering provide an example, with the basic option
  consisting of a 12 square metre parklet with bike repair facilitates,
  and bike parking and/or bike share docking \citep{CoMoUK:0aa}.}. In
\emph{Mobility hub delivery models} \cite{CoMoUK:2021aa} discuss six
types of mobility hubs: ``large interchanges/city hubs\ldots; transport
corridor, smaller interchanges or linking hubs; business park or new
housing development hubs; suburb/mini-hubs; small market town/village
hubs; and tourism hubs'' along with suggested components for each.
Further types might included the `shared' and `electric' forms of
mobility hub, as in the ROBUST sites, wherein the offered modes are
available for hire and electrically powered respectively. It is unclear,
however, whether there is a need to differentiate between electric hubs
(or `e-hubs') where the electrification involves on-site recharging, as
at ROBUST and for many e-cars, and hubs where electric vehicles are
dockless or otherwise `recharged' through the replacement of battery
packs with those powered up elsewhere.

Guidance from \cite{CoMoUK:0ab} suggests a timeline of 12 to 15 months
to step through consultation, design, stakeholder engagement, tendering
and construction, although this appears to relate to either unpowered
mobility hubs or replacable batteries. This guidance also suggests that
budgets i n the vicinity of 250k to 650k pounds \cite{CoMoUK:0ac}; and
that ``mobility hubs work best when part of a wider sustainable
transport network that provides active travel, public transport and
shared transport options'' \cite{CoMoUK:0aa}. The design problem here,
therefore, appears be to allow creation of mobility hubs that:
incorporate existing modes, facilities and services; minimise cost and
timeframes, such that practitioners and tactical urbanists might build
legitimacy through implementation (c.f.
\cite{Reynolds:2020aa,Reynolds:2021ab}) using minimal resources,
incremental approaches that allow expansion one-hub-at-a-time, and that
are scaleable; and allow for rapid iteration, revision or adjustment to
account for changing circumstances. There are limits, however, to what
can be achieved in the context of this research study, suggesting that
the focus here should be on designing tools that allow for mobility hub
creation without the sorts of `hard' engineering works that would
necessitate design, tendering and construction of new facilities or
infrastructure.\\
Rather, this research defines the design problem to be solved as to
create software tools that might facilitate the creation and rapid
development of user information products to support bottom-up, tactical
urbanist and iterative approaches.

\subsection{Development}\label{development}

Both the Andrew Road and various Bristol mobility hubs feature
pedestal-mounted signage that includes: the hub name; iconographic
representations of hub facilities; directional indications; and two
maps. These maps are, respectively, of the local vicinity (a 5-10 minute
walk) of the hub; and the surrounding area (a 15-20 minute walk). A
starting point for developing software packages to support creation of
further mobility hubs, therefore, would appear to be outputting similar
maps of the local and surrounding areas, composited together into a
image of similar size to the pedestal signage encountered during the
site visits along with site naming and other information.

\begin{figure}

{\centering \includegraphics[width=0.5\linewidth]{graphics/layout_test} 

}

\caption{Mockup of hub signage, produced using developed R code. Source: authors.}\label{fig:mock_up_layout_with_maps_grob_output}
\end{figure}

Figure \ref{fig:mock_up_layout_with_maps_grob_output} shows output
software developed during this research, which uses Open Streets Map
(OSM) data sets and mapping tiles to generate a mobility hub information
panel. Here the software has been used to create output for the Galway
ROBUST e-mobility hub, although as an example the `surrounding area map'
shown is for the Letterkenny ROBUST hub location so as to demonstrate
the generalisability of the research outcomes. Key features of the
developed software include the ability to highlight points of interest:
in Figure \ref{fig:mock_up_layout_with_maps_grob_output} various
religious and medical buildings are highlighted in purple, while the
library is shown in blue and nearby schools in orange. This is largely
similar to the approach used in the maps discovered during the site
visits, but might allow considerable flexibility for practitioners in
choosing between colours, iconography or other symbols. A powerful
aspect of adopting the OSM datasets and an automated approach for
creating maps is that settings can be adjusted by the user and new maps
generated within minutes. For example, the `surrounding area map' shown
in Figure \ref{fig:mock_up_layout_with_maps_grob_output} a slightly
different OSM tile option is selected, and buildings recorded as
``pubs'' (rather than public buildings) are shown in blue.

\begin{figure}

{\centering \includegraphics[width=1\linewidth]{graphics/galway_map_small} 

}

\caption{Immediate hub vicinity mockup. Source: authors.}\label{fig:mock_up_small}
\end{figure}

Figure \ref{fig:mock_up_small} further demonstrates the flexibility of
this approach in presenting an `immediate vicinity' map of the Galway
hub. At this or similarly small scales it might be possible to include
individual hub elements (i.e.~the e-bikes or car share vehicles) so that
users can locate the various services and facilities. This represents an
extension of the mapping approaches seen in the site visits, and might
perhaps remove the need for signage or other physical elements to mark
the location of hub elements (e.g.~the green cubes in Bristol). As well,
at this scale individual elements of the surrounding environment may be
more readily apparent to hub users. In Figure \ref{fig:mock_up_small}
the mapping tiles show individual pedestrian footpaths. However, as
these and other layers can be obtained from OSM these and other features
might similarly be mapped, coloured or otherwise highlighted if desired.

\subsection{Delivery}\label{delivery}

The software developed in conjunction with this research is available
online at \textless TO BE INCLUDED AFTER PEER REVIEW. CODE INCLUDED IN
SUBMISSION\textgreater{}

\subsection{Discussion and
conclusions}\label{discussion-and-conclusions}

This design-based research has developed a basic suite of mapping and
publication software tools that might be used in the context of mobility
hubs. The general format and styling is largely based on elements in
place at various mobility hubs in Bristol, but as discussed above the
flexibility of R might allow users to adjust layouts, styling and other
details as desired. Another area in which the developed approach appears
to provide advantages is due to the use of OSM data. As well as being
freely available, the OSM online platform allows this dataset to be
edited by anyone, and there exists a large community of users who
contribute, edit and otherwise maintain this information. In using this
dataset this research may assist towards generalisability and use of the
study outputs by practitioners in that the OSM database is accessible
and has worldwide coverage. As well, because the OSM dataset is
user-editable it may be possible for those who use the outputs of this
research to make corrections and additions as appropriate, thereby
providing an improved dataset for both the purposes of mobility hub
implementation and for other users or use cases. For example, updating
data in OSM for the building immediately north-east of the Westside
Public Library (Figure \ref{fig:mock_up_small}) might allow that
building's name to be displayed on a mobility-hub-related map, yet also
add to the completeness of the OSM dataset overall.

Clearly, there are many avenues for further research in this area. The
software might be expanded to automate the production of websites,
posters, banners, leaflets and other informational products that might
be generated quickly by practitioners, adovocates or others, and so
easily maintained and updated. There might be opportunities to develop
more generic materials for other purposes, for example: maintenance
schedules, inspection logs, management resources and other items that
might help a mobility hub `product owner' facilitate operations,
continuance and growth. Further development of `best practices' in
mobility hub informational provision might involve additional case
studies, especially of non-UK sites, as well as improvements to the base
`hub signage mockup' developed here (e.g.~additional text,
implementation of symbolic mapping for facilities such as toilets,
groceries and cafes, etc.). Furthermore, there may be a need to better
understand what facilities and services might best be included as `part'
of a hub, as opposed to those that are mapped as being `nearby'. More
immediately, user testing and further development of the software, and
of the software outputs, especially the maps, would seem an important
direction for further research efforts in the short term.

This paper's title asked a question related to the applicability of
tactical urbanist approaches to the topic of mobility hubs. Underlying
this may be larger questions related to `what is a mobility hub' and
`who gets to define and advertise something as being a mobility hub'.
Perhaps every railway station, bus or tram stop, shared or
micro-mobility facilities or many other places are \emph{already}
mobility hubs, but just no one has put up a sign or handed out a leaflet
that designates this as being the case? There may be laws and other
impediments limiting the ability of many people to adopt a DIY attitude
to designating, advertising or informing others about how somewhere
might already be, or have the potential to be, a mobility hub. Where
these are less of an issue, however, it is suggested that with the
software developed through this research and access to plotting devices,
especially those capable of outputing larger size materials (e.g.~A0,
A1), almost anyone might be able to `make' a mobility hub in a suitable
location. Maybe even just some smaller-sized leaflets or similar
materials handed or post-out to a local community might help towards
shifting travel behaviours towards more sustainable modes, and encourage
those who can to make more use of what mobility services are already
offered at a local level.

\begin{longtable}[]{@{}llllllll@{}}
\caption{Author contribution matrix, as per CRediT (Contribution Roles
Taxonomy)}\tabularnewline
\toprule\noalign{}
Component & NA & BC & BM & SM & JR & AS & AV \\
\midrule\noalign{}
\endfirsthead
\toprule\noalign{}
Component & NA & BC & BM & SM & JR & AS & AV \\
\midrule\noalign{}
\endhead
\bottomrule\noalign{}
\endlastfoot
Conceptualization & & & & & X & & \\
Data Curation & & & & & X & & \\
Formal Analysis & & & & & X & & \\
Funding Acquisition & & X & & & & X & \\
Investigation & & & & & X & & \\
Methodology & & & & & X & & \\
Project Administration & & & & X & & & \\
Resources & & & & & & & \\
Software & & & & & X & & \\
Supervision & & & X & & & & X \\
Validation & & & & & X & & \\
Visualization & & & & & X & & \\
Writing -- Original Draft Preparation & & & & & X & & \\
Writing -- Review \& Editing & & & X & & X & & \\
\end{longtable}

\renewcommand\refname{References}
\bibliography{References.bib, packages.bib}


\end{document}
